% Section for Interlude E: why 20 regular tetrahedra do not make a regular icosahedron.
% Written to follow W. Kusner's argument. Labels prefixed TET:.

\section{Twenty tetrahedra do not make an icosahedron}
\label{TET:sec}

Frank's icosahedral cluster (and the jitterbug of Lecture~11) rests on a small geometric fact: a regular
icosahedron is \emph{not} made of $20$ regular tetrahedra. Cone the icosahedron from its centre and
you get $20$ congruent tetrahedra, but they are slightly flattened. The edges from the centre are a
little \emph{shorter} than the edges of the icosahedron:
\[
  \frac{\text{circumradius}}{\text{edge}}=\sin\frac{2\pi}{5}\approx0.951<1 .
\]
In terms of balls: twelve unit balls placed at the vertices of an icosahedron around a central unit
ball can all touch the central ball, but then they do not touch each other. This slack is why the
twelve outer balls can roll around (the kissing configuration in $\R^3$ is not rigid), and it is the
root of the Newton--Gregory dispute over whether a thirteenth ball fits.

The usual proof is a computation, for instance with the dihedral angle $\arccos\frac13$ of the regular
tetrahedron. The following argument needs no computation. It uses only the plane fact that \emph{six}
equal circles fit exactly around a seventh, together with the triangle inequality on the sphere.

\begin{proposition}\label{TET:prop}
A regular icosahedron is not the union of $20$ regular tetrahedra sharing a vertex at its centre.
\end{proposition}

\paragraph{The argument.} Suppose it were, with edge length $a$. Put balls of radius $a/2$ at the centre
and at the $12$ vertices. Then the central ball touches all twelve, and each vertex ball touches its
five neighbours. Project the twelve contact points radially onto a sphere about the centre, and
measure distances by angle. Neighbouring vertices are then exactly $60^\circ$ apart, because the centre
and two neighbours form an equilateral triangle.

Stand the icosahedron on a vertex: a north pole $N$, the antipodal vertex $S$, and two
\emph{tropical rings} of five vertices each. The upper ring consists of the neighbours of $N$, the
lower ring of the neighbours of $S$, and the rings are staggered by $36^\circ$ of longitude. The equatorial
band between the rings is made of ten triangles (Figure~\ref{TET:fig:band}, left). Every vertex of the
upper ring is $60^\circ$ from $N$, so the upper ring lies at latitude $30^\circ$\,N. Likewise the lower ring
lies at latitude $30^\circ$\,S.

Now \emph{rotate the lower ring} by $36^\circ$ about the axis $NS$, keeping its latitude. Each lower vertex
$\ell$ now sits directly below an upper vertex $u$, and the ten equatorial triangles have become five
equatorial quadrilaterals, the ``squares'' (Figure~\ref{TET:fig:band}, right). Each of the five
meridians through $N$ and $S$ now carries four of our points,
\[
  N,\ u,\ \ell,\ S\qquad\text{at angles}\qquad 0^\circ,\ 60^\circ,\ 120^\circ,\ 180^\circ
\]
along the meridian, since $u$ is $60^\circ$ from $N$, $\ell$ is $60^\circ$ from $S$, and $N$ and $S$ are
$180^\circ$ apart. The meridian is a great circle, and $60^\circ$ steps are exactly the spacing of a
\emph{ring of six} balls around the central one, as in the hexagonal packing
(Figure~\ref{TET:fig:meridian}). So after the rotation, $u$ and $\ell$ touch: the five meridians have
become five half-rings of six, one great circle each.

But before the rotation $u$ and $\ell$ also touched, since they were neighbours in the icosahedron, and
in between we only slid $\ell$ sideways along its latitude circle. The closest point of a latitude
circle to $u$ is the one on $u$'s meridian, so sliding $\ell$ off the meridian moves it strictly
\emph{farther} from $u$. In other words, the deformation from ten equatorial triangles to five
equatorial squares, made without changing any latitude, shortened the distance from $u$ to $\ell$. Yet
it is $60^\circ$ at both ends. That is a contradiction. \qed

\medskip
In one line, the same argument is the triangle inequality on the sphere:
\[
  d(u,\ell)\ \ge\ d(N,\ell)-d(N,u)\ =\ (180^\circ-60^\circ)-60^\circ\ =\ 60^\circ,
\]
with equality only if $u$ lies on the arc from $N$ to $\ell$, that is, on $\ell$'s meridian. Each upper
vertex $u$ has two lower neighbours, and they cannot both lie on $u$'s meridian, so for at least one of
them $d(u,\ell)>60^\circ$. (In the icosahedron both are off the meridian, by $\pm36^\circ$.) The ring of six
is the picture of the equality case.

\begin{figure}[ht]
\centering
\begin{tikzpicture}[scale=0.95]
  % left: antiprism band (10 triangles), right: prism band (5 squares)
  \begin{scope}
    \foreach \k in {0,...,5} {\coordinate (U\k) at (\k,1); }
    \foreach \k in {0,...,4} {\coordinate (L\k) at (\k+0.5,0); }
    \coordinate (Lm) at (-0.5,0);
    \draw[thick] (U0)--(U5); \draw[thick] (Lm)--(L4);
    \foreach \k [evaluate=\k as \j using int(\k+1)] in {0,...,4} {\draw (U\k)--(L\k)--(U\j);}
    \draw (U0)--(Lm);
    \foreach \k in {0,...,5} \fill (U\k) circle (0.06);
    \foreach \k in {0,...,4} \fill[flagcol] (L\k) circle (0.06);
    \node[font=\small] at (2.5,1.55) {upper ring, latitude $30^\circ$\,N};
    \node[font=\small] at (2.5,-0.55) {lower ring, latitude $30^\circ$\,S};
    \node[font=\small,gray] at (2.5,-1.1) {ten triangles (icosahedron)};
  \end{scope}
  \draw[-{Stealth},thick] (5.8,0.5)--(6.9,0.5) node[midway,above,font=\scriptsize]{rotate};
  \node[font=\scriptsize] at (6.35,0.2) {lower ring};
  \begin{scope}[xshift=7.4cm]
    \foreach \k in {0,...,5} {\coordinate (U\k) at (\k,1); \coordinate (L\k) at (\k,0);}
    \draw[thick] (U0)--(U5); \draw[thick] (L0)--(L5);
    \foreach \k in {0,...,5} \draw (U\k)--(L\k);
    \foreach \k in {0,...,5} {\fill (U\k) circle (0.06); \fill[flagcol] (L\k) circle (0.06);}
    \node[font=\small,gray] at (2.5,-1.1) {five squares, same latitudes};
  \end{scope}
\end{tikzpicture}
\caption{The equatorial band, unrolled (the left and right ends are identified; longitudes in units of
$72^\circ$). Rotating the lower ring by $36^\circ$ turns the ten equatorial triangles into five
quadrilaterals without changing any latitude.}
\label{TET:fig:band}
\end{figure}

\begin{figure}[ht]
\centering
\begin{tikzpicture}[scale=0.62]
  \fill[gray!15] (0,0) circle (1); \draw (0,0) circle (1); \node[font=\scriptsize] at (0,0) {$O$};
  \foreach \a/\lab in {90/N,30/u,-30/\ell,-90/S}
    {\filldraw[fill=addcol!15,draw=addcol,thick] (\a:2) circle (1); \node[font=\small] at (\a:2) {$\lab$};}
  \foreach \a in {150,210} {\draw[dashed,gray] (\a:2) circle (1);}
  \draw[gray,thin] (0,0) circle (2);
  \foreach \a in {90,30,-30,-90} \draw[gray,thin] (0,0)--(\a:2);
  \node[font=\scriptsize] at (60:0.65) {$60^\circ$};
  \node[font=\scriptsize] at (0:0.7) {$60^\circ$};
  \node[font=\scriptsize] at (-60:0.65) {$60^\circ$};
\end{tikzpicture}
\caption{After the rotation, each meridian is half of a ring of six balls around the central ball: $N$,
$u$, $\ell$, $S$ are consecutive at $60^\circ$ steps. The dashed balls complete the ring.}
\label{TET:fig:meridian}
\end{figure}

\begin{remark}[What the argument shows, quantitatively]
In the hypothetical configuration with rings at latitudes $\pm30^\circ$, an upper and a lower vertex
whose longitudes differ by $36^\circ$ are at angular distance $\arccos\big(\tfrac34\cos36^\circ-\tfrac14\big)
\approx69.1^\circ$, not $60^\circ$. In the true icosahedron the rings sit lower, at latitudes
$\pm\arctan\frac12\approx\pm26.6^\circ$, every edge subtends $\arctan2\approx63.4^\circ$ at the centre, and
the gaps between neighbouring outer balls are what let the twelve balls roll.
\end{remark}

\begin{remark}[Two related classical facts]
(1) Regular tetrahedra do not even fill the space around an \emph{edge}: the dihedral angle is
$\arccos\frac13\approx70.53^\circ$, and five such angles leave a gap of about $7.36^\circ$. (Aristotle
claimed that regular tetrahedra tile space; see Lagarias--Zong, \emph{Mysteries in packing regular
tetrahedra}, Notices AMS 59 (2012).) The same angle appears here: five points at $60^\circ$ from $N$ and
$60^\circ$ from each other would need angles of $\arccos\frac13$ at $N$, and five of these do not close up.
(2) In curved space the frustration disappears. In the $3$-sphere, $600$ regular tetrahedra (with
the right curvature) fit together exactly, five around each edge, forming the regular polytope
$\{3,3,5\}$ with $120$ vertices, each surrounded by a regular icosahedron. This is the ``ideal''
icosahedral order of Frank's liquids; the flat-space packing problem has to compromise.
\end{remark}

\begin{remark}[Source]
\unclear{Wording of the attribution is a placeholder for the author to edit.}
This argument came out of a conversation between the author and T.~C.~Hales.
\end{remark}
